\section{Introduction}

\begin{figure}[h]
    \centering
    \includegraphics[width=0.95\columnwidth]{figure/Intro.pdf} 
    \caption{An example comparing generic and personalized query expansion.}
    % \caption{Same query, personalized needs.\textcolor{red}{the title needs modification}}
    \label{fig:intro}
    \vspace{-20pt}
\end{figure}
Retrieval-Augmented Generation (RAG) has emerged as a pivotal paradigm for enhancing the capabilities of large language models (LLMs)~\cite{lewis2020retrieval, zhao2024retrieval,zhang2025lsrp}, leveraging a two-stage process: first retrieving relevant external corpus, then generating responses conditioned on the retrieved information~\cite{cheng2023lift}. 
The effectiveness of this paradigm hinges on the quality of query expansion by using LLM world knowledge~\cite{culpepper2021topic,song2024survey}, which directly impacts the accuracy and relevance of the retrieval stage. 
Most current strategies generate pseudo answers~\cite{gao2023precise} or candidate completions~\cite{jia2023mill} using LLMs in a zero-shot or few-shot manner. 
These generated semantic representations, often combined with the original query, are used to retrieve information from a global external corpus.
% Retrieval-Augmented Generation (RAG) has emerged as a pivotal paradigm for enhancing the capabilities of large language models (LLMs)~\cite{lewis2020retrieval, zhao2024retrieval,zhang2025lsrp}, leveraging a two-stage process: first retrieving relevant external corpus, then generating responses conditioned on the retrieved information~\cite{cheng2023lift}. 
% The effectiveness of this paradigm hinges on the quality of query expansion by using LLM world knowledge~\cite{culpepper2021topic,song2024survey}, which directly impacts the accuracy of the retrieval stage. 
% Most current strategies generate pseudo answers~\cite{gao2023precise} or candidate completions~\cite{jia2023mill} using LLMs in a zero-shot or few-shot manner. 
% These generated semantic representations, often combined with the original query, are used to retrieve information from a global external corpus.

% Retrieval-Augmented Generation (RAG) has emerged as a pivotal paradigm for enhancing the performance of large language models (LLMs)~\cite{lewis2020retrieval, zhao2024retrieval,zhang2025lsrp}, leveraging a two-stage process that first retrieves external knowledge and then generates responses conditioned on the retrieved content~\cite{cheng2023lift}. The quality of query formulation~\cite{culpepper2021topic,song2024survey} during the retrieval stage critically determines the relevance of retrieved documents and, consequently, the accuracy of the generated outputs. Existing approaches typically improve retrieval by rewriting or expanding queries to better capture semantic intent. However, these methods often overlook the individual semantic characteristics of users. In real-world applications, the same surface-level query may convey vastly different intentions depending on a user's preferences~\cite{zou2023users}, background knowledge~\cite{westhofen2025temporal}, and expression style~\cite{neelakanteswara2024rags}. This lack of personalized semantics in query understanding limits the generalization and effectiveness of current RAG systems in user-facing scenarios.

% Despite recent progress, existing query expansion strategies~\cite{jia2023mill,gao2023precise,wang2023query2doc} remain limited in their ability to capture user-specific semantics.
% \yy{Most methods generate pseudo answers or candidate completions using LLMs without conditioning on user-specific histories, treating all queries with the same expansion procedure regardless of individual linguistic or contextual preferences.}
% As illustrated in Figure~\ref{fig:intro} (top), when both users submit the same query, ``\emph{How would you describe my dietary preferences?}'', existing methods generate an identical expansion. This leads to mismatched retrieval: while the health-conscious user (User 1) receives content emphasizing ``\emph{organic ingredients}'' and balanced nutrition, the variety-seeking user (User 2)—whose preferences involve ``\emph{trying new side dishes}'' and casual experimentation—retrieves irrelevant results. This failure stems from ignoring valuable cues embedded in user history, such as linguistic tendencies, habitual intent framing, or thematic focus. Without incorporating these personalized signals, expanded queries remain semantically misaligned with user intent, limiting the retrieval performance.
% 具体来说一句，传统的方法做了什么

% Many existing approaches focus on expanding queries with LLM world knowledge. 
% Some methods enrich queries by generating semantically informed completions~\cite{wang2023query2doc}, while others iteratively refine the query using generative feedback loops~\cite{mackie2023generative}. 
% Most current strategies generate pseudo answers~\cite{gao2023precise} or candidate completions~\cite{jia2023mill} using LLMs, applying a uniform query expansion procedure. 
% These methods are evaluated in zero-shot or few-shot settings, assuming retrieval over a global, externally defined corpus with non-personalized queries.
% These techniques are often evaluated in zero-shot or few-shot settings and assume retrieval in a general corpus and related query. 

Recently, with LLMs widely adopted as personalized assistants, many user queries are required to retrieve from individual histories and contexts, demanding retrieval mechanisms that go beyond generic corpora to leverage user-specific data.
Existing query expansion strategies remain limited in their ability to capture user-specific semantics. 
In practice, the same textual query may convey vastly different intentions depending on a user’s preferences~\cite{zou2023users}, background knowledge~\cite{westhofen2025temporal}, or expression style~\cite{neelakanteswara2024rags}. As users exhibit diverse personas and contextual needs, it becomes increasingly important for query expansion methods to be personalized—capable of adapting to user-specific semantics within private corpora—so as to ensure accurate retrieval and enhance downstream generation.
% Existing expansion strategies remain limited in their ability to capture user-specific semantics. In real-world scenarios, the same user query in text may reflect vastly different intentions depending on a user’s preferences~\cite{zou2023users}, background knowledge~\cite{westhofen2025temporal}, or expression style~\cite{neelakanteswara2024rags}. This lack of personalized understanding during query expansion significantly restricts the generalization and effectiveness of retrieval in the user corpus. 
% As illustrated in Figure~\ref{fig:intro} (top), when two users submit the same query—``\emph{How would you describe my dietary preferences?}''—standard expansion methods produce generic outputs. This leads to mismatched retrieval: for example, the health-conscious user (User 1) is shown content related to ``\emph{organic ingredients},'' are successful retrievaled. While the variety-seeking user (User 2), whose personal content involve ``\emph{trying new side dishes},'' could not be retrievaled. 

As illustrated in Figure~\ref{fig:intro} (top), when two users submit the same query—``\emph{How would you describe my dietary preferences?}''—standard expansion methods yield generic outputs that fail to account for individual differences. This results in mismatched retrieval: for example, the health-conscious user (User~1) successfully retrieves content related to ``\emph{organic ingredients},'' which aligns with their preferences. In contrast, the variety-seeking user (User~2), whose personal content emphasizes ``\emph{trying new side dishes},'' fails to retrieve relevant information, highlighting the limitations of non-personalized expansion strategies. Such failures stem from neglecting rich signals in user—such as style, intent patterns, and thematic focus—resulting in semantically misaligned expansions and degraded retrieval performance.
% Such failures arise from ignoring rich cues embedded in user history, including stylistic tendencies, recurring intent structures, and thematic focus. Without incorporating these signals, expanded queries remain semantically misaligned with user intent, ultimately degrading retrieval performance.


Motivated by these limitations, we aim to develop a personalized query expansion framework that adapts to individual user intent and expression as shown in Figure~\ref{fig:intro}. However, introducing personalization into expansion brings two fundamental challenges:
\textbf{(1) User expression styles are inherently diverse.}  
Users articulate intent using varied linguistic patterns—ranging from minimal prompts to elaborative reasoning—driven by personal habits, domain familiarity, or rhetorical preference~\cite{neelakanteswara2024rags}. These styles are often implicit and non-transferable, making it difficult to construct expansions that faithfully preserve user-specific semantics within a generalizable framework.
\textbf{(2) User corpora induce heterogeneous semantic structures.} 
Users’ personal corpora often differ significantly in topical coverage, content organization, and linguistic granularity. This high degree of heterogeneity makes it difficult to locate reliable semantic regions related to the user query, and in turn, complicates the alignment of user-specific preferences. Without user-specific semantic anchors to guide expansion, queries are prone to drifting toward irrelevant regions, resulting in mismatched retrieval despite plausible surface semantics.
These challenges raise a central question: \textit{How can we personalize query expansion—both in style and structure—before retrieval to align with individual user intent and corpus characteristics?}
% Personalized corpora vary significantly in topical density, lexical scope, and organizational granularity. This structural heterogeneity not only fragments the semantic embedding space but also obscures the underlying patterns in user expression. Without adaptive anchoring mechanisms, it becomes increasingly difficult to expand queries that remain semantically anchored and stylistically faithful.

% Personalized query expansion, while promising, faces two fundamental challenges that undermine its effectiveness in real-world settings. 
% \yy{While query expansion has demonstrated utility in improving retrieval quality, existing methods face significant limitations when applied to personalized settings. We identify two challenges:}
% \textbf{(1) Standard expansions fail to preserve user-specific expression styles.}  
% Users often communicate intent through distinct linguistic patterns—ranging from minimal prompts to elaborate descriptions—shaped by personal habits, domains, or preferences~\cite{neelakanteswara2024rags}. These stylistic variations are not superficial: they directly influence how intent is encoded and interpreted. Generic expansion models, trained on aggregated data or fixed templates, fail to account for such diversity, leading to semantic drift between the expanded query and the user’s true intent.
% \textbf{(2) Uniform strategies struggle with diverse user corpora structures.}  
% User-specific corpora often exhibit non-uniform topical focus, lexical granularity, and structural coherence. Such heterogeneity disrupts the consistency of semantic representations, making it difficult for expanded queries to anchor meaningfully within the embedding space induced by general-purpose retrievers. Without personalized anchoring, even semantically plausible expansions may retrieve contextually irrelevant results.
% These challenges raise a central question: \textit{How can we personalize query expansion—both in style and structure—before retrieval to align with individual user intent and corpus characteristics?}


To address these challenges, we propose a \textit{Personalize Before Retrieve} (\textbf{PBR}) framework that integrates user-specific signals into query understanding prior to retrieval. PBR adopts a two-stage design that combines expression-level expansion with structure-level alignment, enabling queries to reflect both personal style and corpus semantics.
\textbf{\ding{182} P-PRF: Personal Style-Aligned Pseudo Relevance Feedback.}  
This module extracts linguistic patterns from user history to guide LLMs in generating pseudo-utterances and reasoning paths. These signals capture personalized expression style and hidden intent, often missed by generic expansion methods.
\textbf{\ding{183} P-Anchor: Personal Structure-Aligned Semantic Anchoring.}  
We encode user history into a semantic graph that captures localized relevance patterns. Personalized PageRank is applied to identify representative anchor points within this structure to encode the whole user corpora or to express user general preference. The query is then guided toward these anchors to reflect the user’s semantic landscape.
By jointly modeling stylistic variation and structural relevance, P-PRF and P-Anchor produce personalized query representations that enhance alignment between user intent and retrieval semantics in RAG systems. Our contributions can be summarized as follows:
\begin{itemize}[leftmargin=*]
    \item To the best of our knowledge, we are the first to propose a personalized query expansion framework for RAG systems that adapts to the user for improving retrieval.
    \item We propose \textbf{P-PRF}, a style-aware expansion module that generates pseudo queries and reasoning paths conditioned on user history.
    \item We introduce \textbf{P-Anchor}, a graph-based alignment module that grounds queries within personalized semantic spaces via local corpus structure.
    \item Experiments on two real-world dialogue datasets show that \textbf{PBR} significantly improves personalized retrieval, validating the benefits of pre-retrieval personalization.
\end{itemize}


% Personalized query expansion, while promising, is hindered by two fundamental challenges that we identify in this work. 
% \textbf{1) User-Specific Expression Styles Break General Query Expansion.} Users often exhibit distinct linguistic styles—some favor~\cite{neelakanteswara2024rags} concise formulations while others prefer elaborative and richly descriptive expressions. Even when posing identical information needs, the underlying intent may manifest through different syntactic structures, lexical choices, or rhetorical patterns. This stylistic variability renders template-based or globally trained expansion models inadequate, as they fail to accurately capture individual expressive habits, leading to semantic drift between the expanded query and the user's actual intent.
% \textbf{2) Structural Heterogeneity in User Corpora Breaks Semantic Anchoring.} User-specific corpora are inherently heterogeneous in terms of topical focus, structural organization, and granularity of expression. As a result, queries often struggle to locate appropriate semantic anchors~\cite{bogdan2025thought} within a shared embedding space. Without personalized guidance rooted in the user's content distribution, the expansion process risks drifting into semantically irrelevant regions, undermining the relevance of subsequent retrieval results.
% These observations motivate our central research question: \textbf{\textit{How can we model user-specific expression styles and align with users' own data prior to retrieval, in order to construct personalized query representations tailored for retrieval within individualized corpora?}}
